%! Operations with Polynomials — Unit 3, Lesson 3.1 — Guided Notes (Answer Key)
\documentclass[10pt]{article}
\usepackage{algebra2-article}
\usepackage{algebra2-key}   % pulls in -boxes; do NOT also load -boxes
\newcommand{\TallMath}[1]{$\displaystyle #1\rule[-1.4em]{0pt}{3.2em}$}

\begin{document}

\pageheader{Unit 3, Lesson 3.1}{Guided Notes --- Answer Key}
\namedateperiod

\vspace{0.08in}

\begin{objectivebox}
\textbf{By the end of this lesson, I will be able to\dots}
\begin{itemize}\setlength{\itemsep}{2pt}
  \item write a polynomial in \emph{standard form} and name its \emph{degree}, \emph{leading coefficient}, and number of \emph{terms};
  \item \emph{add} and \emph{subtract} polynomials by combining like terms, distributing a subtraction to \emph{every} term;
  \item \emph{multiply} polynomials in one and two variables and apply the \emph{special products} $(a\pm b)^2$ and $(a+b)(a-b)$;
  \item predict a product's \emph{degree} and \emph{leading coefficient} and state the \emph{end behavior} they force.
\end{itemize}
\end{objectivebox}

\vspace{0.05in}

\begin{vocabbox}
\small
Fill in each term as we build it together.
\vspace{2pt}
\termblanklong{Standard form}
\ansline{Terms written from the highest power down to the constant: $a_nx^n+\dots+a_1x+a_0$.}
\termblanklong{Like terms}
\ansline{Terms with identical variable parts --- same variables, same exponents. Only like terms combine.}
\termblanklong{Monomial / binomial / trinomial}
\ansline{A polynomial with one, two, or three terms.}
\termblanklong{Degree of a product}
\ansline{$\deg(fg) = \deg f + \deg g$ --- degrees add, and the leading coefficients multiply.}
\termblanklong{Special products}
\ansline{$(a+b)^2 = a^2+2ab+b^2$; \; $(a-b)^2 = a^2-2ab+b^2$; \; $(a+b)(a-b) = a^2-b^2$.}
\end{vocabbox}

\begin{hookbox}
Yesterday we used \emph{two} names for the same curve:
\[
  g(x) = (x-1)^2(x+2) \qquad\text{and}\qquad g(x) = x^3 - 3x + 2 .
\]
Check a few inputs in \emph{both} forms:
\begin{center}
\renewcommand{\arraystretch}{1.3}
\begin{tabular}{|c|c|c|}
\hline
$x$ & $(x-1)^2(x+2)$ & $x^3-3x+2$ \\ \hline
$0$ & \ans{2} & \ans{2} \\ \hline
$2$ & \ans{4} & \ans{4} \\ \hline
$-2$ & \ans{0} & \ans{0} \\ \hline
\end{tabular}
\end{center}
\begin{itemize}\setlength{\itemsep}{1pt}
  \item Agreeing at three points is good evidence. Is it \emph{proof} that the two forms are the same function? Why or why not?
        \ansline{No --- many different curves can pass through the same three points. Matching at finitely many inputs never proves two expressions agree everywhere.}
  \item What would actually prove it?
        \ansline{Multiplying the factored form out and showing it becomes $x^3-3x+2$ term for term.}
\end{itemize}
Each costume is useful: the \textbf{factored} form hands you the \emph{zeros}; the \textbf{standard}
form hands you the \emph{degree} and \emph{leading coefficient}, hence the \emph{end behavior}. Today
we learn to travel from the first to the second.
\end{hookbox}

\begin{notesbox}{1. Naming a Polynomial}
Always put the polynomial in \textbf{standard form} first --- powers from \ans{highest} to
\ans{lowest} --- because every other name depends on it.

\vspace{3pt}
Rewrite $-5x^2 + 7x^4 - 3 + x$ in standard form: \ans{$7x^4 - 5x^2 + x - 3$}

\vspace{3pt}
\begin{itemize}[leftmargin=1.5em, itemsep=2pt]
  \item \textbf{Degree:} \ans{4} \qquad \textbf{Leading coefficient:} \ans{7} \qquad
        \textbf{Constant term:} \ans{$-3$}
  \item \textbf{Number of terms:} \ans{4}, so it is a (monomial / binomial / trinomial / none of
        these): \ans{none of these}
  \item In \emph{two} variables, the degree of a \emph{term} is the \textbf{sum} of its exponents. So
        $-7x^2y^3$ has degree \ans{5}.
\end{itemize}
\end{notesbox}

\begin{notesbox}{2. Adding \& Subtracting: Collecting, Not Creating}
\textbf{Like terms} have the \emph{same} variables raised to the \emph{same} exponents. So $3x^2$ and
$3x^3$ \ans{do not} combine.

\vspace{4pt}
\textbf{Add.} $(3x^3 - 5x + 2) + (x^3 + 4x^2 - 7) = $ \ans{$4x^3 + 4x^2 - 5x - 5$}

\vspace{5pt}
\textbf{Subtract --- the whole lesson in one move.} A minus sign in front of a parenthesis is a
factor of $-1$, so it changes the sign of \ans{every} term inside.

\vspace{3pt}
$(3x^3 - 5x + 2) - (x^3 + 4x^2 - 7)$

\vspace{2pt}
\emph{Step 1 --- rewrite with every sign changed:} \; $3x^3 - 5x + 2$ \ans{$-\,x^3 - 4x^2 + 7$}

\vspace{2pt}
\emph{Step 2 --- collect like terms:} \ans{$2x^3 - 4x^2 - 5x + 9$}

\vspace{5pt}
\textbf{A degree surprise.} $(x^4 + 2x) - (x^4 - 5) = $ \ans{$2x + 5$}, which has degree
\ans{1} --- not $4$! The leading terms \ans{cancel}, so a sum or difference has degree
\emph{at most} the larger of the two degrees.
\end{notesbox}

\begin{notesbox}{3. Multiplying: Every Term Times Every Term}
One rule --- the \textbf{distributive property} --- in three sizes. Remember: coefficients
\ans{multiply}, exponents \ans{add}.

\vspace{4pt}
\textbf{(a) Monomial $\times$ polynomial.} \; $-2x^2(3x^3 - 4x + 6) = $ \ans{$-6x^5 + 8x^3 - 12x^2$}

\vspace{5pt}
\textbf{(b) Binomial $\times$ binomial} (four products --- ``FOIL'' is just a name for the same
distribution). \; $(2x - 5)(3x + 4)$

\vspace{2pt}
Four products: \ans{$6x^2$} $+$ \ans{$8x$} $+$ \ans{$(-15x)$} $+$ \ans{$(-20)$}

\vspace{2pt}
Collected: \ans{$6x^2 - 7x - 20$}

\vspace{5pt}
\textbf{(c) Binomial $\times$ trinomial} (six products --- use the box so nothing is dropped).
\; $(x + 3)(x^2 - 2x + 4)$
\begin{center}
\renewcommand{\arraystretch}{1.9}
\begin{tabular}{|c|c|c|c|}
\hline
$\times$ & $x^2$ & $-2x$ & $+4$ \\ \hline
$x$   & \ans{$x^3$} & \ans{$-2x^2$} & \ans{$4x$} \\ \hline
$+3$  & \ans{$3x^2$} & \ans{$-6x$} & \ans{$12$} \\ \hline
\end{tabular}
\end{center}
Collect like terms: \ans{$x^3 + x^2 - 2x + 12$}

\vspace{5pt}
\textbf{(d) Two variables.} \; $(3x + 2y)(x - 4y) = $ \ans{$3x^2 - 10xy - 8y^2$}

\vspace{2pt}
The two middle products, $-12xy$ and $+2xy$, are like terms because they have the
\ans{same variable part, $xy$}.

\vspace{5pt}
\textbf{(e) The shape of a product --- before you expand.} In $(2x^3 - 1)(-3x^4 + x)$, only the two
\emph{leading} terms can make the highest power: $2x^3 \cdot (-3x^4) = $ \ans{$-6x^7$}.
\begin{itemize}[leftmargin=1.5em, itemsep=2pt]
  \item Degrees \ans{add}: \; $3 + 4 = $ \ans{7}. \quad Leading coefficients
        \ans{multiply}: \; $2 \cdot (-3) = $ \ans{$-6$}.
  \item So (Lesson 3.0) the degree is \emph{odd} and the leading coefficient is \emph{negative}:
        the left arm goes \ans{up} and the right arm goes \ans{down}.
\end{itemize}
\end{notesbox}

\begin{notesbox}{4. Special Products --- Worth Knowing on Sight}
Derive them once, then recognize them.
\begin{itemize}[leftmargin=1.5em, itemsep=3pt]
  \item \textbf{Square of a sum:} $(a+b)^2 = (a+b)(a+b) = $ \ans{$a^2 + 2ab + b^2$}
  \item \textbf{Square of a difference:} $(a-b)^2 = $ \ans{$a^2 - 2ab + b^2$}
  \item \textbf{Difference of squares:} $(a+b)(a-b) = $ \ans{$a^2 - b^2$} \; --- the two middle terms
        \ans{cancel}.
\end{itemize}

\vspace{4pt}
\textbf{The trap.} $(x+4)^2$ is \emph{not} $x^2 + 16$. Test it at $x = 1$: \;
$(1+4)^2 = $ \ans{25} \; but \; $1^2 + 16 = $ \ans{17}. \; The missing piece is the
\textbf{middle term} \ans{$8x$}. Correct answer: \ans{$x^2 + 8x + 16$}

\vspace{4pt}
Now use the identities:
\begin{itemize}[leftmargin=1.5em, itemsep=3pt]
  \item $(3x - 4)^2 = $ \ans{$9x^2 - 24x + 16$}
  \item $(2x + 7)(2x - 7) = $ \ans{$4x^2 - 49$}
  \item $(5x - 3y)^2 = $ \ans{$25x^2 - 30xy + 9y^2$}
\end{itemize}
\end{notesbox}

\begin{practicebox}
\textbf{Back to the Hook.} Prove that $(x-1)^2(x+2)$ really is $x^3 - 3x + 2$.
\begin{enumerate}[leftmargin=1.7em, itemsep=6pt]
  \item First square the binomial: $(x-1)^2 = $ \ans{$x^2 - 2x + 1$} \; (special product --- don't forget the
        middle term).
  \item Now multiply that trinomial by $(x+2)$. Six products:
        \begin{center}
        \renewcommand{\arraystretch}{1.9}
        \begin{tabular}{|c|c|c|c|}
        \hline
        $\times$ & $x^2$ & $-2x$ & $+1$ \\ \hline
        $x$   & \ans{$x^3$} & \ans{$-2x^2$} & \ans{$x$} \\ \hline
        $+2$  & \ans{$2x^2$} & \ans{$-4x$} & \ans{$2$} \\ \hline
        \end{tabular}
        \end{center}
  \item Collect like terms and write the result in standard form: \ans{$x^3 - 3x + 2$}
  \item Does it match $x^3 - 3x + 2$? \ans{Yes} \; Its degree is \ans{3} and its leading
        coefficient is \ans{1}, so the arms go \ans{left down, right up} --- exactly the graph you read in
        Lesson 3.0.
\end{enumerate}
\end{practicebox}

\begin{teachernote}
The $-2x^2$ and $+2x^2$ cells in the practice box are the payoff --- they cancel, which is precisely
why the standard form has \emph{no} $x^2$ term. Point at that cancellation explicitly; students who
expect every box cell to survive are the ones who later mis-collect. In section~2, the ``degree
surprise'' is worth a full minute: it is the only place all year where a degree drops, and it makes
the contrast with multiplication (where degrees always add) land. If time is short, section~4's
two-variable item $(5x-3y)^2$ can move to homework, but do \emph{not} cut the $x=1$ numeric test of
$(x+4)^2$ --- it is the cheapest permanent fix for that error.
\end{teachernote}

\end{document}
